\chapter{Liposuction}

\section*{What Is This Surgery?}
Liposuction is an elective surgical procedure designed for body contouring, specifically targeting and removing localized, stubborn deposits of subcutaneous adipose tissue that are resistant to conventional methods of diet and exercise. It is primarily an aesthetic intervention, not a weight-loss solution, aimed at improving the shape, proportion, and overall silhouette of specific body areas. The technique involves the insertion of a thin, hollow tube, known as a cannula, through small, strategically placed incisions into the fatty layer beneath the skin. This cannula is then meticulously manipulated to mechanically disrupt and emulsify fat cells, which are subsequently aspirated using a vacuum pump or large syringe. Commonly performed on regions such as the abdomen, flanks, thighs, hips, buttocks, arms, neck, and chin, liposuction offers a refined and more sculpted appearance, enhancing self-image and complementing a healthy lifestyle.

\section*{Alternative Names}
\begin{itemize}
  \item Lipoplasty
  \item Suction-Assisted Lipectomy (SAL)
  \item Suction Lipectomy
  \item Fat Aspiration
  \item Body Contouring Surgery
  \item Lipo-sculpture
\end{itemize}

\section*{Relevant Surgical Anatomy}
Understanding the intricate anatomy of the subcutaneous tissue is paramount for safe and effective liposuction. The skin is composed of the epidermis and dermis, beneath which lies the subcutaneous fat layer, also known as the hypodermis or panniculus adiposus. This layer is organized into superficial and deep compartments by a fibrous fascial system. The superficial fat layer, located immediately beneath the dermis, is typically dense and fibrous, containing a rich network of neurovascular structures. The deep fat layer, situated above the muscle fascia, is generally looser and more amenable to fat removal. Key fascial planes, such as Scarpa's fascia in the abdomen and Camper's fascia, delineate these compartments and are crucial surgical landmarks.

Arterial supply and venous drainage to the subcutaneous fat originate from perforating vessels that traverse the deeper muscle fascia. These vessels, along with cutaneous nerves, run within the fibrous septa that compartmentalize the fat. Major neurovascular bundles, particularly in areas like the neck, axilla, and medial thigh, must be meticulously avoided to prevent injury. Lymphatic drainage generally follows the venous system, and excessive trauma can disrupt these channels, leading to prolonged edema. Surgical landmarks, such as the iliac crests, umbilicus, and inframammary fold, guide the surgeon in achieving symmetrical and harmonious contours while respecting underlying vital structures.

\section{Surgical Principle \& Rationale}
The fundamental surgical principle of liposuction revolves around the selective mechanical disruption and aspiration of adipose tissue from targeted subcutaneous regions, while preserving the overlying skin and underlying vital structures. The rationale is to permanently reduce the number of fat cells in specific areas, thereby reshaping the body contour. The procedure typically commences with the infiltration of a tumescent solution, a dilute mixture of crystalloid saline, lidocaine (a local anesthetic), and epinephrine (a vasoconstrictor), into the fatty layer. This solution serves a triple purpose: it provides profound local anesthesia, significantly constricts blood vessels to minimize intraoperative bleeding and postoperative bruising, and engorges the fat cells, making them easier to dislodge and aspirate.

Once the tumescent solution has taken effect, small, inconspicuous incisions are created. A specialized cannula, connected to a vacuum source, is then inserted into the fat layer. The surgeon employs a controlled, fanning motion to mechanically dislodge adipocytes and create a network of tunnels within the adipose tissue. The dislodged fat cells, suspended in the tumescent fluid, are then suctioned out. The technical goals are to achieve a smooth, even reduction of fat, create harmonious body contours, and promote skin retraction over the newly sculpted areas, all while minimizing trauma to surrounding tissues and avoiding irregularities. The permanent removal of fat cells in the treated area means that these specific cells cannot re-expand, though remaining fat cells elsewhere can, underscoring the importance of stable weight maintenance for long-term results.

\section{History \& Surgical Evolution}
The concept of fat removal for aesthetic purposes has its nascent roots in early 20th-century dermatological procedures, though these early attempts were often crude and associated with significant complications, including skin necrosis and infection, due to aggressive techniques and a limited understanding of tissue planes. Modern liposuction techniques began to emerge and gain traction in the 1970s.

A pivotal breakthrough occurred in 1977 when French surgeons Arpad and Giorgio Fischer developed the blunt cannula technique. This innovation significantly minimized trauma to blood vessels and nerves compared to sharp curettage, marking a crucial step towards safer fat removal. However, it was Dr. Yves-G\'{e}rard Illouz, also from France, who is widely credited with popularizing and refining the technique in the early 1980s. Illouz introduced the "wet technique," involving the injection of saline into the fatty tissue to facilitate fat removal and reduce blood loss, making the procedure more manageable and safer. Further advancements in the 1980s, particularly by American dermatologist Dr. Jeffrey Klein, led to the development of the "tumescent technique." This revolutionary method involved the infiltration of large volumes of highly dilute lidocaine and epinephrine, which dramatically improved safety by reducing blood loss to negligible levels and allowed for the procedure to be performed under local anesthesia, establishing it as the gold standard. Subsequent technological evolutions include ultrasound-assisted liposuction (UAL), power-assisted liposuction (PAL), and laser-assisted liposuction (LAL), which utilize energy or mechanical vibration to further emulsify fat, making aspiration easier and potentially more efficient, though the core tumescent principle remains central to most modern approaches.

\section{Clinical Indications}
Liposuction is indicated for individuals seeking to improve their body contours and address specific areas of localized fat accumulation that are resistant to diet and exercise. Common clinical indications include:
\begin{itemize}
  \item To reduce disproportionate fat deposits in the abdomen and flanks (often referred to as "love handles") that create an undesirable silhouette.
  \item To sculpt the outer thighs ("saddlebags"), inner thighs, and knees to achieve a more streamlined and aesthetically pleasing leg contour.
  \item To diminish excess fat on the hips and buttocks, creating a smoother, more balanced, and harmonious posterior profile.
  \item To refine the upper arms, addressing "bat wings" and improving arm definition and tone.
  \item To remove submental fat, commonly known as a "double chin," and to enhance the definition of the jawline and neck area.
  \item To treat pseudogynecomastia in men, which is the presence of fatty tissue in the male breast area, distinct from true glandular gynecomastia.
  \item To improve overall body proportion and to serve as an adjunct or refinement procedure following other body contouring surgeries, such as abdominoplasty.
  \item To address localized fat accumulation associated with lipedema, a chronic condition characterized by abnormal fat distribution, though this requires specialized expertise.
\end{itemize}

\section*{Clinical Positioning}
Liposuction is almost exclusively positioned as a primary, elective cosmetic surgical procedure. It is considered a first-line surgical option for individuals who are at or near their ideal body weight but struggle with localized fat deposits that do not respond to conventional weight management strategies. It is crucial to emphasize that liposuction is not intended as a primary treatment for obesity, nor is it a substitute for a healthy diet and regular exercise regimen. Its role is to refine and sculpt, not to achieve significant weight loss. In certain clinical scenarios, liposuction may be considered a secondary or adjunctive procedure. For instance, it can be used to refine the results of a previous body contouring surgery, such as an abdominoplasty (tummy tuck), by addressing residual fat pockets. It may also be employed to contour specific areas after significant weight loss, including post-bariatric surgery patients, to improve overall body shape. However, its most common and primary role remains as a standalone procedure for aesthetic enhancement in otherwise healthy individuals with realistic expectations.

\section{Surgical Technique \& Procedure}
The liposuction procedure is meticulously performed following a structured sequence to ensure patient safety and optimal aesthetic outcomes.
\begin{itemize}
  \item \textbf{Anesthesia:} The choice of anesthesia is tailored to the patient and the extent of the procedure. Options include local anesthesia with conscious sedation for smaller areas, regional anesthesia (e.g., epidural) for larger regions, or general anesthesia for extensive cases or patient preference. This decision is made during the preoperative consultation with the surgeon and anesthesiologist.
  \item \textbf{Patient Positioning and Marking:} The patient is positioned to allow optimal access to all target areas. While standing, the surgeon meticulously marks the areas of fat accumulation, highlighting both areas for removal and transition zones, along with the planned incision sites. This standing assessment is critical as gravity alters fat distribution when lying down.
  \item \textbf{Tumescent Infiltration:} Small, discreet incisions (typically 3-5 mm) are made. A sterile infiltration cannula is then used to evenly distribute the tumescent solution (saline, lidocaine, epinephrine) throughout the subcutaneous fat layer of the marked areas. Adequate time (15-30 minutes) is allowed for the solution to take full effect, ensuring vasoconstriction and anesthesia.
  \item \textbf{Fat Aspiration:} After tumescence, a different, often slightly larger, aspiration cannula is inserted through the same incisions. The surgeon systematically moves the cannula in a controlled, back-and-forth, fanning motion through the fatty tissue. This mechanical action dislodges adipocytes and creates a network of tunnels. The dislodged fat and tumescent fluid are then suctioned out via a vacuum pump or large syringe. The surgeon continuously palpates and assesses the contour to ensure smooth, even fat removal and avoid irregularities.
  \item \textbf{Contour Assessment and Symmetry:} Throughout the aspiration phase, the surgeon frequently pauses to visually and tactilely assess the evolving body contour, ensuring symmetry and smoothness across all treated areas. In some cases, the patient may be briefly repositioned to allow for a comprehensive assessment from multiple angles.
  \item \textbf{Closure:} Once the desired amount of fat has been removed and the contour is satisfactory, the small incisions are typically left open or closed with one or two fine sutures. Leaving them slightly open can facilitate the drainage of residual tumescent fluid, which helps to reduce postoperative swelling and bruising.
  \item \textbf{Compression Garment Application:} Immediately following the procedure, a specialized compression garment is applied to the treated areas. This garment is crucial for reducing swelling, minimizing fluid accumulation (seroma formation), promoting skin retraction, and encouraging the skin to adhere smoothly to the new underlying contours. No drains or implants are typically used in standard liposuction.
\end{itemize}

\section*{Preoperative Workup \& Preparation}
Thorough preoperative workup and preparation are paramount for ensuring patient safety and optimizing the outcomes of liposuction surgery.
A comprehensive medical history and physical examination are conducted to assess the patient's overall health status and identify any potential risks. This typically includes:
\begin{itemize}
  \item \textbf{Laboratory Tests:} Routine blood tests such as a complete blood count (CBC), coagulation profile (PT/INR, PTT), basic metabolic panel, and sometimes a pregnancy test for women of childbearing age.
  \item \textbf{Cardiac Evaluation:} For older patients or those with pre-existing cardiovascular conditions, an electrocardiogram (ECG) and potentially a cardiac clearance from a cardiologist may be required.
  \item \textbf{Medication Review:} Patients must provide a complete list of all medications, supplements, and herbal remedies. Blood-thinning agents (e.g., aspirin, NSAIDs, warfarin, clopidogrel, certain herbal supplements like ginkgo biloba, garlic, or vitamin E) must typically be discontinued for one to two weeks prior to surgery to minimize the risk of excessive bleeding.
  \item \textbf{Smoking Cessation:} Patients are strongly advised to cease smoking at least 4-6 weeks before surgery, as nicotine impairs wound healing and significantly increases the risk of complications.
  \item \textbf{NPO Status:} Patients undergoing general anesthesia or significant sedation are instructed to refrain from eating or drinking (NPO) for a specified period, usually 6-8 hours, before surgery to prevent aspiration.
  \item \textbf{Prophylactic Antibiotics:} A single dose of prophylactic antibiotics may be prescribed to be taken orally before or intravenously at the time of the procedure to reduce the risk of surgical site infection.
  \item \textbf{Informed Consent:} A detailed discussion of the procedure, its potential benefits, risks, alternatives, and expected outcomes is conducted, and the patient provides fully informed consent. Establishing realistic expectations is a critical component of this process.
  \item \textbf{Pre-operative Photography:} Standardized photographs of the treatment areas are taken from multiple angles for documentation, surgical planning, and objective comparison of results.
\end{itemize}

\section*{Contraindications \& Precautions}
\textbf{Absolute Contraindications:}
\begin{itemize}
  \item Severe systemic diseases, including uncontrolled diabetes mellitus, severe cardiovascular disease (e.g., unstable angina, recent myocardial infarction), significant pulmonary disease, or uncompensated kidney or liver failure, which substantially increase surgical and anesthetic risk.
  \item Active infection in the proposed treatment area or a systemic infection (sepsis).
  \item Diagnosed coagulation disorders or patients on chronic anticoagulation therapy that cannot be safely discontinued for the perioperative period.
  \item Pregnancy or breastfeeding.
  \item Unrealistic patient expectations regarding the surgical outcome or a diagnosis of body dysmorphic disorder, which may indicate psychological unsuitability for elective cosmetic surgery.
  \item Documented allergy to lidocaine, epinephrine, or other components of the tumescent solution.
\end{itemize}

\textbf{Relative Contraindications and Precautions:}
\begin{itemize}
  \item Significant obesity (Body Mass Index (BMI) greater than 30-35 kg/m\textasciicircum2): Liposuction is not a weight-loss procedure, and patients who are significantly overweight may achieve suboptimal aesthetic results and face higher complication rates.
  \item Poor skin elasticity or significant skin laxity: May result in loose, sagging skin after fat removal, potentially necessitating additional skin excision procedures (e.g., abdominoplasty, brachioplasty) to achieve a smooth contour.
  \item History of deep vein thrombosis (DVT) or pulmonary embolism (PE): Requires careful risk assessment, thromboprophylaxis, and potentially consultation with a hematologist.
  \item Certain medications: Immunosuppressants, high-dose steroids, or medications known to impair wound healing may require adjustment or discontinuation.
  \item Presence of hernias in the treatment area: These should ideally be repaired prior to or concurrently with liposuction to prevent complications and ensure a smooth contour.
  \item Previous abdominal surgery or other surgeries in the treatment area: May increase the risk of adhesions, making cannula passage more challenging and potentially increasing the risk of visceral perforation.
  \item Age: While not an absolute contraindication, older patients may have reduced skin elasticity and slower healing capabilities, which can impact the final aesthetic outcome.
\end{itemize}

\section*{Complications \& Risks}
While generally considered safe when performed by experienced, board-certified surgeons, liposuction, like any surgical procedure, carries potential risks and complications. Patients should be thoroughly informed of these possibilities.
\begin{itemize}
  \item \textbf{General Surgical Complications:}
    \begin{itemize}
      \item \textit{Bleeding:} Hematoma formation (collection of blood under the skin) or excessive intraoperative blood loss, particularly with large-volume liposuction.
      \item \textit{Infection:} Localized surgical site infection, cellulitis, or, rarely, more severe systemic infections requiring antibiotic treatment or surgical drainage.
      \item \textit{Anesthesia Complications:} Adverse reactions to anesthetic agents, respiratory depression, cardiac arrhythmias, or other systemic events related to general or regional anesthesia.
      \item \textit{Fluid Imbalance and Electrolyte Disturbances:} Especially with large-volume liposuction and extensive tumescent fluid infiltration, there is a risk of fluid overload, pulmonary edema, or electrolyte abnormalities (e.g., hyponatremia).
    \end{itemize}
  \item \textbf{Procedure-Specific Risks:}
    \begin{itemize}
      \item \textit{Contour Irregularities:} Uneven fat removal, skin dimpling, waviness, depressions, or asymmetry, which may necessitate revision surgery to correct.
      \item \textit{Seroma:} Accumulation of clear, yellowish fluid under the skin, sometimes requiring aspiration or drainage.
      \item \textit{Numbness or Altered Sensation:} Temporary or, rarely, permanent changes in skin sensation (hypoesthesia or hyperesthesia) due to nerve irritation or damage during cannula manipulation.
      \item \textit{Skin Discoloration:} Persistent bruising, hyperpigmentation (darkening of the skin), or hypopigmentation (lightening of the skin) in the treated areas.
      \item \textit{Skin Necrosis:} A rare but serious complication where skin tissue dies due to compromised blood supply, potentially requiring wound care or skin grafting.
      \item \textit{Fat Embolism:} An extremely rare but life-threatening complication where fat particles enter the bloodstream and travel to the lungs (pulmonary fat embolism) or brain, leading to respiratory distress or neurological deficits.
      \item \textit{Deep Vein Thrombosis (DVT) and Pulmonary Embolism (PE):} Formation of blood clots in the deep veins of the legs, which can dislodge and travel to the lungs, a serious and potentially fatal complication. Prophylactic measures are crucial.
      \item \textit{Scarring:} While incisions are small, visible scars can occur, particularly in individuals prone to hypertrophic or keloid scarring.
      \item \textit{Thermal Injury:} Possible with energy-assisted liposuction techniques (e.g., UAL, laser-assisted liposuction) if not performed carefully, leading to burns to the skin or deeper tissues.
      \item \textit{Perforation of Internal Organs:} Extremely rare, but possible if the cannula penetrates too deeply, particularly in the abdomen (bowel perforation) or chest cavity.
    \end{itemize}
\end{itemize}

\section*{Postoperative Recovery Timeline}
Immediately after liposuction, patients are transferred to a Post-Anesthesia Care Unit (PACU) for close monitoring as they emerge from anesthesia. Pain, swelling, and bruising are universally expected. A compression garment is applied to the treated areas in the operating room and is crucial for recovery, typically worn continuously for several weeks as instructed by the surgeon. Drainage of residual tumescent fluid from the small incision sites is common for the first 24-48 hours and is a normal, beneficial part of the healing process, helping to reduce swelling and bruising. Patients are usually discharged home the same day, provided they are stable and have a responsible adult caregiver.

During the first 1-2 weeks post-procedure, patients should anticipate significant swelling, bruising, and some discomfort, which is usually manageable with prescribed oral pain medication. Light activity, such as short walks, is strongly encouraged from day one to promote circulation and minimize the risk of deep vein thrombosis, but strenuous exercise, heavy lifting, and activities that strain the treated areas should be strictly avoided. Most patients can return to light, non-strenuous work or daily activities within a few days to a week, depending on the extent of the procedure and their individual healing capacity. Long-term recovery involves the gradual resolution of swelling, which can take several months, with the final aesthetic results typically becoming apparent between 3 to 6 months post-surgery. Maintaining a stable weight through a healthy diet and regular exercise is paramount for preserving the long-term aesthetic outcome.

\section*{Surgical Findings \& Pathology}
During liposuction, the primary surgical findings documented by the surgeon include the specific anatomical areas treated, the type and volume of tumescent solution infiltrated, and the total volume of aspirate removed. The aspirate itself is a mixture of fat cells, tumescent fluid, and a small amount of blood. The surgeon visually assesses the quality and character of the aspirate, noting its color and consistency, which can provide immediate feedback on the effectiveness of fat removal and the degree of blood loss. The goal is to achieve a smooth, even contour with appropriate fat reduction, which is assessed by palpation and visual inspection throughout the procedure.

Formal pathological examination of the aspirated fat tissue is not routinely performed for standard cosmetic liposuction. The resected tissue is adipose tissue, and its macroscopic appearance is typically sufficient for surgical purposes. However, in rare instances where a suspicious mass or lesion is encountered during the procedure, or if there is a clinical concern for an underlying pathological process (e.g., liposarcoma, atypical lipomatous tumor), a sample of the tissue may be sent for histopathological analysis. The pathology report, in such cases, would describe the cellular composition, presence of any atypical cells, and confirm the diagnosis, guiding any further management. For the vast majority of liposuction cases, the "pathology" is simply confirmation of benign adipose tissue.

\section{Expected Postoperative Course}
A normal and successful postoperative course following liposuction is characterized by a predictable progression of healing and resolution of expected symptoms, culminating in an improved body contour. Initially, patients will experience moderate pain, which typically subsides significantly within the first few days and can be managed with oral analgesics. Swelling and bruising are most pronounced during the first 1-2 weeks and then gradually diminish over the subsequent weeks to months. The compression garment plays a crucial role in mitigating these symptoms and promoting skin retraction. Patients should expect a gradual return to normal activities, with light walking encouraged early on, and a progressive increase in physical activity over 4-6 weeks. The skin in the treated areas will initially feel firm or indurated due to swelling and tissue remodeling, but this sensation typically softens over several months. The small incision sites will heal, leaving inconspicuous scars that fade over time. The ultimate goal is for the treated areas to exhibit a smoother, more sculpted appearance with a noticeable reduction in localized fat deposits, achieving symmetry and a harmonious contour without visible irregularities. Patient satisfaction, coupled with the objective improvement in body shape, signifies a successful outcome.

\section{Postoperative Care \& Follow-up}
Comprehensive postoperative care and diligent follow-up are essential for monitoring healing, managing potential complications, and ensuring optimal long-term results after liposuction.
\begin{itemize}
  \item \textbf{Initial Follow-up Visits:} Patients typically have their first follow-up appointment within 5-7 days post-surgery to assess wound healing, remove any sutures if present, and evaluate the initial reduction in swelling. Subsequent visits are usually scheduled at 2-4 weeks, 3 months, and 6-12 months to monitor progress, assess contour, and address any concerns.
  \item \textbf{Compression Garment Wear:} Continued and consistent wear of the prescribed compression garment is crucial for several weeks (typically 4-6 weeks full-time, then part-time as advised by the surgeon) to reduce swelling, minimize seroma formation, and promote optimal skin retraction and adherence to the new contours.
  \item \textbf{Activity Resumption:} Patients are guided on a gradual return to normal activities. Light walking is encouraged immediately, but strenuous exercise, heavy lifting, and activities that put pressure on the treated areas are restricted for 4-6 weeks.
  \item \textbf{Massage and Lymphatic Drainage:} Some surgeons may recommend manual lymphatic drainage massages or self-massage techniques to help reduce persistent swelling, improve circulation, and soften any areas of firmness or induration.
  \item \textbf{Weight Management Counseling:} Emphasis is placed on maintaining a healthy lifestyle, including a balanced diet and regular exercise, to preserve the long-term results of the surgery, as weight gain can diminish the aesthetic outcome.
  \item \textbf{Warning Signs:} Patients are thoroughly educated on warning signs that necessitate immediate contact with their surgeon, including fever (above 101$^{\circ}$F or 38.3$^{\circ}$C), increasing redness, warmth, or pus-like discharge from incision sites (signs of infection), excessive or sudden pain, sudden swelling, shortness of breath, or calf pain (potential signs of DVT/PE).
\end{itemize}
Routine imaging or laboratory tests are not typically required after uncomplicated liposuction, unless specific complications are suspected.

\section*{Epidemiology}
Liposuction consistently ranks among the most frequently performed cosmetic surgical procedures globally, reflecting its widespread demand for body contouring. According to statistics from major aesthetic surgery societies, hundreds of thousands of liposuction procedures are performed annually worldwide, often placing it among the top three cosmetic surgical interventions, alongside breast augmentation and rhinoplasty. The demographic profile of liposuction patients is broad, encompassing both men and women, typically ranging in age from their early 20s to their 60s, with a peak incidence observed in individuals in their 30s and 40s. While women historically constitute a larger proportion of liposuction patients, there has been a notable and steady increase in male liposuction procedures, particularly targeting areas such as the abdomen, flanks, and chest (for pseudogynecomastia). The most common anatomical sites for liposuction include the abdomen, flanks, thighs, and arms. Mortality rates associated with liposuction are exceedingly low, estimated to be less than 1 in 50,000 procedures, especially when performed by board-certified plastic surgeons in accredited facilities and adhering to safe volume limits (typically less than 5 liters of pure fat aspirate). Morbidity, primarily related to minor complications such as contour irregularities, seroma formation, and temporary numbness, is more common but generally manageable and rarely life-threatening.

\section*{Outcomes \& Prognosis}
The outcomes of liposuction are generally favorable, with high rates of patient satisfaction when realistic expectations are established preoperatively. Typical success rates for achieving significant and noticeable contour improvement range from 85\% to 95\%. Functional outcomes are primarily aesthetic, leading to enhanced body image, improved clothing fit, and a significant boost in self-confidence and psychological well-being. Survival outcomes are not directly impacted by liposuction, as it is an elective cosmetic procedure performed on generally healthy individuals. The long-term prognosis for maintaining the achieved results is excellent, provided the patient adheres to a stable weight. Fat cells removed during liposuction are permanently excised and do not regenerate; however, the remaining fat cells in treated areas, or fat cells in untreated areas, can still enlarge if weight is gained, potentially diminishing the aesthetic result. Recurrence of localized fat deposits in the *treated* areas is rare if weight is maintained, but new fat accumulation can occur in other body regions if a healthy lifestyle is not sustained. Prognostic factors for optimal outcomes include good inherent skin elasticity, a stable pre-operative weight close to the ideal, absence of significant medical comorbidities, and strict adherence to postoperative care instructions, including consistent compression garment wear. Patients with poor skin elasticity may experience residual skin laxity post-liposuction, which might necessitate secondary procedures like skin excision (e.g., abdominoplasty, thigh lift) to achieve a tauter contour.

\section*{References}
\begin{enumerate}
  \item American Society of Plastic Surgeons. Liposuction. Available at: \url{https://www.plasticsurgery.org/cosmetic-procedures/liposuction}. Accessed February 2025.
  \item Rohrich RJ, Beran SJ. Liposuction. Quality Medical Publishing; 2024.
  \item Matarasso A. Liposuction. In: Thorne CH, Chung KC, Gosain AK, Gurtner GC, Mehrara BJ, Rubin JP, eds. Grabb and Smith's Plastic Surgery. 8th ed. Wolters Kluwer; 2024.
  \item MedlinePlus. Liposuction. U.S. National Library of Medicine; 2025.
\end{enumerate}

\clearpage